% 【理论扩展】所属：第3章（由 chapters/revised/03-justification-context.tex \input 引入）
\minitopic{知识归属的语义证据}若“知道”具有语境敏感性，它应在语言行为中表现出某些规律。典型比较对象是“高”：说某人高，要看比较群体。语境主义者认为，“知道”也隐含一组待排除可能性；谈话目的和显著反例改变这组范围\parencite{derose2009}。普通机场接人时，“我知道航班九点到”只需正常查询；事故调查中，同一句话可能要求核实系统时区、更新延迟和临时改降。 反对者指出，“高”允许自然的程度修饰和明确比较：“对小学生来说很高”。“知道”却不自然地说“对低标准而言知道”。我们还会在跨语境转述时坚持：“昨天我说我知道，今天没有发现新证据，所以我仍知道。”这似乎把知识当作主体稳定状态，而非归属者语境函数。 语境主义可回应，许多语义敏感项并不共享全部表面语法；知识标准也可由隐含相关替代决定。真正检验需看否认、回指、间接引语和撤回行为，而非仅凭一个银行案例。语义理论应解释一组语言现象，不能只配合预先偏好的认识论结论。

\minitopic{怀疑论如何改变显著性}Lewis提出，知识要求证据排除一切没有被适当忽略的可能性；某些可能性可忽略，实际情形、与证据相似的情形等则不能忽略\parencite{lewis1996}。提及一个错误可能往往使它不能继续被忽略，这解释了怀疑论的对话力量。 但“提到就相关”会让知识过度脆弱。一个人恶意喊出“也许你的手机是一块完美道具”，不应自动撤销所有权知识。显著性需要与证据、会话任务或公认规则结合。纯粹想象与存在具体错误机制之间，应有认识差别。 我们还要区分心理显著与规范相关。飞机事故刚上新闻，会使人高估飞行风险；显著性提高并不表示风险在证据上更相关。语境理论若把任何心理突出都变成知识标准，可能把偏差写进语义；若只接受规范显著，又需提供非语义的相关性理论。

\minitopic{不变主义的替代解释}严格不变主义认为，“知道”的标准在所有语境相同，日常判断差异来自语用。低风险中说“我知道”传达足够把握；高风险中即使字面为真，这样说也可能暗示不需复核，因而不恰当。听者把语用不当误判为字面为假。 温和不变主义则设定一个不要求绝对确定的固定门槛。它解释跨语境稳定性，却需说明银行案例为何强烈变化。语用、注意和行动规范可以共同承担解释。理论比较的关键不是谁能讲一个故事，而是谁用更独立、统一的原则解释更多数据。 主体敏感不变主义承认知识地位变化，但变化来自主体实践情形而非词义\parencite{stanley2005}。高风险使主体需要更强证据。它保留“知道”的固定语义，却使同一证据在不同利害下产生不同知识，面对知识是否应纯粹认识性的质疑。

\minitopic{纯粹认识理由与实践理由}纯粹主义主张，决定信念认识合理性的只是真理相关因素，如证据和可靠性；金钱、愿望和后果不改变证据。实践主义则认为，何种信心适合行动可进入知识或确证。争论不能靠“事实不会随利益改变”解决，因为双方都同意真值不变；问题是知识门槛是否只由真值关系决定。 设两人有相同检验结果，一人是否立即手术关系生命，另一人仅决定是否在课堂举例。实践主义可能要求前者复核。纯粹主义可说，两人对诊断拥有相同认识地位，但行动理由不同。要区分两者，可以询问前者在不涉及行动的考试中是否仍知道诊断；若直觉恢复，变化可能主要是行动规范。 实践理由还可能影响调查义务而非当前确证。高风险主体有义务收集更多证据，但在收集前仍可能知道。类似地，消防员知道楼内有人，却仍须二次确认房间位置。把“应继续查”直接推出“不知道”需要额外原则。

\minitopic{知识行动规范的细化}“只以所知为行动前提”过强：紧急撤离可以基于尚未达到知识的火警；“只要高概率就可行动”又忽略错误代价和证据类型。决策理论通常把概率与效用结合，知识规范则强调某些命题可被视作固定背景。二者可能分工：概率决定风险权衡，知识决定哪些前提无需在当前决策中重新打开。 实际行动包含可逆性。发送一封可撤回提醒与实施不可逆手术，对证据要求不同。实践侵入理论若只看损失大小，可能遗漏收集信息成本、决策期限和纠错机会。更完整的实践情形包括错误后果、行动可逆性、进一步证据价值和责任角色。 断言规范也有强弱。直接无保留断言可能要求知识；明确概率、来源或保留语气的报告只需相应较弱地位。一个气象员说“模型给出70\%降雨概率”可在不知道“明天会下雨”时完全恰当。良好认识交流通过语气匹配证据，而非迫使所有内容达到同一门槛。

\minitopic{案例组：同证据是否同知识}\textbf{旅行案例：}两人都看见护照在桌上。一人明早出国，另一人只是整理房间。高风险是否使前者不知道护照在桌上，还是只要求他把护照收入包中？由于知觉极直接，很多人认为知识不变，支持行动义务解释。 \textbf{邮件案例：}两人都记得邮件截止日期。一人迟交损失巨大，另一人无后果。记忆可错且复核成本低，知识判断更容易随风险变化。差异可能来自读者把低成本复核当作应尽义务，而非风险直接进入知识。 \textbf{数学案例：}两人都正确心算 $37\times 19$，高风险者决定航天参数。即使他知道答案，制度仍要求机器与同事复核。这里“知道”显然不足以免除程序。案例说明，实践需要冗余不等于个人没有知识。

\minitopic{综合案例：法庭证据门槛}刑事与民事审判使用不同证明标准，但这是否表示同一陪审员在刑事案中不知道、在民事案中知道同一事实？法律标准决定国家可采取何种行动，不必直接定义个体知识。把“排除合理怀疑”当作知识定义，会把制度价值和错误分配混入日常概念。 另一方面，法律语境确实改变相关替代和调查义务。证据链缺口在闲谈中可忽略，在剥夺自由的决定中不可忽略。最稳妥分析是分层：事实证据给予某种置信和可能知识；法律规范规定行动门槛；语境决定哪些质疑必须公开回答。三层可互动，但不应彼此替代。

\begin{tcolorbox}[breakable,colback=white,colframe=blue!30!black,title={语境案例控制表},fonttitle=\bfseries]
确保成对案例具有相同证据、同一主体能力、同一时间；分别改变归属者谈话、主体利害或具体反证；记录判断变化是在“知道”、应当复核、可否断言还是可否行动；再比较语境主义、不变主义与实践侵入的解释成本。
\end{tcolorbox}
